In this work, we investigated how differences in model capability, game complexity, and competition shape both individual and collective outcomes in strategic interactions between LLM agents.
Our preliminary findings indicate a non-monotonic relationship between adversary capability and baseline utility, as well as a `skewed deal' phenomenon in highly competitive settings.

There are several important directions for future work.
One is to generalize our negotiation environments further to larger numbers of agents and items, and to richer preference structures with complementarities or contingent value (for example, only valuing one item if it is accompanied by another).
Our preliminary three-agent extension of Game~1 (Section~\ref{sec:game1_n3}) already suggests that genuinely multi-agent strategic behavior does not simply dilute a capability advantage: with a single strong focal model against two weaker baselines, the focal model's utility edge grew rather than shrank, with the stronger agent acting as a Schelling point around which the weaker agents failed to coordinate a blocking coalition.
However, this extension is based on one replicate per condition and only a single Game~1 competition level, so it should be treated as a preliminary signal rather than a tested scaling result.
A natural next step is therefore to run a fully powered $N > 2$ sweep across all three games, systematically varying both the number and the composition of agents (for example, multiple strong models against a single weak one, or mixed-capability pools), so that we can test whether the Schelling-point pattern we observed is robust and whether there are regimes in which weaker models can in fact coordinate or form coalitions that partially offset the advantage of a single stronger model.
Such settings would also let us test whether the relationship between Arena Elo and strategic performance remains monotonic once coalition formation, agenda-setting, and coordination costs become central.

Another direction is to expand the range of LLM agents we study, including reasoning models whose effective strength can be modulated by limits on their reasoning-token budgets.
This would enable a cleaner investigation of test-time compute scaling in strategic settings, including whether additional reasoning consistently improves bargaining and coordination performance, or whether there are diminishing returns---and even strategic regressions---in some regimes.

Finally, it would be valuable to extend this line of work beyond negotiation to other strategic domains and AI paradigms, such as MARL.
